\documentclass[11pt]{article}
\usepackage[margin=1in]{geometry}
\usepackage[T1]{fontenc}
\usepackage[utf8]{inputenc}
\usepackage{textcomp}
\usepackage{lmodern}
\usepackage{microtype}
\usepackage{xcolor}
\usepackage{amsmath,amssymb,amsthm,mathtools}
\usepackage{enumitem}
\usepackage{setspace}
\usepackage{xurl}
\usepackage[hidelinks]{hyperref}
\providecommand{\tightlist}{\setlength{\itemsep}{0pt}\setlength{\parskip}{0pt}}
\setstretch{1.18}
\setlength{\parindent}{0pt}
\setlength{\parskip}{6pt plus 2pt minus 1pt}
\setlist{itemsep=2pt,topsep=4pt}
\setcounter{secnumdepth}{-\maxdimen}
\emergencystretch=3em
\sloppy
\begin{document}
\section*{The finiteness theorem for the invariants of finite groups}

Let the finite group $G$ consist of the $h$ linear transformations (with nonvanishing determinant)
\[
A_1,\dots,A_h,
\]
where $A_k$ is represented by the linear transformation
\[
x^{(k)}_\nu=\sum_{\mu=1}^n a^{(k)}_{\nu\mu}x_\mu
\qquad (\nu=1,\dots,n),
\]
or, more briefly, by
\[
(x^{(k)})=A_k(x).
\]
Thus the group $G$ carries the row $(x)$ with elements $x_1,\dots,x_n$ into the rows $(x^{(k)})$ with elements $x^{(k)}_1,\dots,x^{(k)}_n$. Since the identity must occur among $A_1,\dots,A_h$, the row $(x)$ itself also occurs among the rows $(x^{(k)})$.

By an \emph{integral rational (absolute) invariant} of the group I mean an integral rational function of $x_1,\dots,x_n$ that remains identically unchanged under application of $A_1,\dots,A_h$; thus for such an invariant $f(x)$ one has
\[
f(x)=f(x^{(1)})=f(x^{(2)})=\cdots=f(x^{(h)}).
\tag{1}
\]

\subsection*{1. A full invariant system from the Galois resolvent}

Formula \textup{(1)} expresses that $f(x)$ is an integral rational symmetric function of the rows $(x^{(1)}),\dots,(x^{(h)})$ --- in fact, since each summand $f(x)$ contains only the single row $(x)$, this is the simplest case, usually called the \emph{uniform} case. By the known theorem on symmetric functions of rows, $f(x)$ can therefore be expressed integrally and rationally through the elementary symmetric functions of these rows, that is, through the coefficients $G_{\alpha_1\dots\alpha_n}(x)$ of the ``Galois resolvent''
\[
\Phi(z,u)=\prod_{k=1}^h\bigl(z-(u_1x^{(k)}_1+u_2x^{(k)}_2+\cdots+u_nx^{(k)}_n)\bigr)
      = z^h+\sum G_{\alpha_1\dots\alpha_n}(x)\,u_1^{\alpha_1}\cdots u_n^{\alpha_n} z^{\alpha_0},
\tag{2}
\]
where the $G_{\alpha_1\dots\alpha_n}(x)$ are invariants of degree $\alpha_1+\cdots+\alpha_n$ in the variables $x$.

Thus we have proved:

\medskip
\noindent
\textbf{Theorem.}
\emph{The coefficients $G_{\alpha_1\dots\alpha_n}(x)$ of the Galois resolvent form a complete system of invariants of the group, in such a way that every invariant can be expressed integrally and rationally in terms of these finitely many invariants.}

\subsection*{2. A second complete system}

Starting from \textup{(1)} one can also make the following still more elementary observation, which at the same time leads to a second complete system.

Suppose
\[
f(x)=a+b\,x_\nu+\cdots+c\,x_\nu^m+\cdots,
\]
where $a,b,\dots,c$ are constants. Then from \textup{(1)} one has
\[
h\cdot f(x)
 = h\cdot a
   + b\sum_{k=1}^h x^{(k)}_\nu
   + \cdots
   + c\sum_{k=1}^h \bigl(x^{(k)}_\nu\bigr)^m
   + \cdots.
\]
Hence every invariant can be assembled integrally and linearly from the special quantities
\[
s_m=\sum_{k=1}^h \bigl(x^{(k)}_\nu\bigr)^m
\qquad (m=0,1,2,\dots),
\]
so it is enough to prove the theorem for these special functions.

Up to a numerical factor, $s_m$ is the coefficient of $u_1^{\alpha_1}\cdots u_n^{\alpha_n}$ in the expression
\[
h\cdot s_m=\sum_{k=1}^h\bigl(u_1x^{(k)}_1+\cdots+u_nx^{(k)}_n\bigr)^m,
\]
where $u=u_1+\cdots+u_n$ is written only as a convenient abbreviation, and this expression represents the $u$-power sum of the $h$ linear forms.

Now it is well known that the infinitely many power sums $S_m$ are integral rational functions of $S_1,S_2,\dots,S_h$; their coefficients are determined by the invariants
\[
S_m=\sum_{k=1}^h\bigl(x^{(k)}_1+\cdots+x^{(k)}_n\bigr)^m
\qquad (m\le h).
\]
Thus we obtain a second complete system:

\medskip
\noindent
\textbf{Theorem.}
\emph{A complete invariant system of the group is given by all invariants $J_{\alpha_1\dots\alpha_n}$ for which}
\[
\alpha_1+\cdots+\alpha_n\le h,
\]
\emph{where $h$ is the order of the group.}

The connection with part~1 is given by the remark that the power sums $S_1,\dots,S_h$ may be replaced by the elementary symmetric functions of $x^{(k)}_1,\dots,x^{(k)}_n$, which leads back to the complete system there given by the $G_{\alpha_1\dots\alpha_n}(x)$.

Both results show that all invariants can be expressed integrally and rationally through those whose degree in the $x$ does not exceed the order $h$ of the group.

\subsection*{3. Rational representation}

From what has just been proved one obtains consequences for rational representation. Every rational absolute invariant can --- as is seen immediately by multiplying numerator and denominator by the conjugates of the denominator with respect to the group --- be represented as the quotient of two integral rational absolute invariants, not necessarily relatively prime. It follows that every rational absolute invariant can be expressed rationally through the coefficients $G_{\alpha_1\dots\alpha_n}(x)$ of the Galois resolvent $\Phi(z,u)$.

This theorem can also be proved in various easy ways without passing through the finiteness theorem; it is already formulated in Weber II, \S 58.

\subsection*{4. Final remark}

Finally, I should mention that the results derived here are already implicitly contained in my paper \emph{Körper und Systeme rationaler Funktionen}: namely, the complete system given in part~1 appears there in Theorems VI and VII, while an argument for rational representation independent of the present finiteness theorem appears there in Theorem III. Both the proof and the result become, however, considerably simpler in the special situation treated here than in the general theory.

I was led to apply those general investigations to the invariants of finite groups by conversations with E.~Fischer, to whom the proof given in part~2 is due in substance as well. In the general case the complete system arising there is not rationally definable, and the remark in part~3 likewise goes back to him.

\begin{flushright}
Erlangen, May 1915.
\end{flushright}

\end{document}
